\begin{figure}[t]
\centering
\begin{tikzpicture}[
    box/.style={
        rectangle,
        draw=black!70,
        fill=black!5,
        minimum width=4.5cm,
        minimum height=1.0cm,
        align=center,
        font=\footnotesize,
        rounded corners=2pt
    },
    layer/.style={
        rectangle,
        draw=blue!50!black,
        fill=blue!5,
        minimum width=5.5cm,
        minimum height=1.4cm,
        align=center,
        font=\footnotesize,
        rounded corners=2pt
    },
    gate/.style={
        rectangle,
        draw=green!50!black,
        fill=green!8,
        minimum width=3.8cm,
        minimum height=1.4cm,
        align=center,
        font=\footnotesize,
        rounded corners=2pt
    },
    decision/.style={
        diamond,
        draw=black!70,
        fill=orange!8,
        minimum width=1.8cm,
        minimum height=1.4cm,
        align=center,
        font=\footnotesize,
        aspect=2
    },
    arrow/.style={->, >=stealth, thick},
    feedback/.style={->, >=stealth, thick, red!60, dashed}
]

% Top: Model Output
\node[box] (model) at (0, 0) {\textbf{Model Output}\\(assistant message)};

% Layer 1
\node[layer, minimum height=2.0cm] (sentinel) at (0, -2.5) {};
\node[font=\footnotesize\bfseries] at (0, -1.85) {Layer 1: Completion Sentinel};
\node[font=\scriptsize, text width=5cm, align=center] at (0, -2.8) {Transcript parsing (jq)\\6 mechanical checks\\Exit 0 / Exit 2};

\draw[arrow] (model) -- (sentinel) node[midway, right, font=\scriptsize] {Stop hook};

% Decision diamond
\node[decision] (decide) at (0, -4.8) {Pass?};
\draw[arrow] (sentinel) -- (decide);

% Block feedback loop
\node[font=\scriptsize, text=red!60] at (-3.8, -3.8) {stderr feedback};
\draw[feedback] (decide.west) -- +(-1.5, 0) |- (model.west)
    node[pos=0.25, below, font=\scriptsize, text=red!60] {Block};

% Layer 2
\node[layer, minimum height=1.6cm] (verifier) at (0, -6.8) {};
\node[font=\footnotesize\bfseries] at (0, -6.3) {Layer 2: File Verifier};
\node[font=\scriptsize, text width=5cm, align=center] at (0, -7.1) {Reads sentinel JSON\\Filesystem rechecks};

\draw[arrow] (decide) -- (verifier) node[midway, right, font=\scriptsize] {Pass};

% Layer 3: VIBE_GATE (side flow)
\node[gate, minimum height=2.0cm] (gate) at (6.5, -2.5) {};
\node[font=\footnotesize\bfseries] at (6.5, -1.85) {Layer 3: VIBE\_GATE};
\node[font=\scriptsize, text width=3.5cm, align=center] at (6.5, -2.85) {35 markers, 5 skills\\Bash command substitution\\values from shell\\Cannot be fabricated};

% Skill execution feeds into VIBE_GATE
\node[box, minimum width=3.5cm] (skill) at (6.5, 0) {\textbf{Skill Execution}\\(during task)};
\draw[arrow] (skill) -- (gate) node[midway, right, font=\scriptsize] {generates};

% VIBE_GATE feeds into Layer 1 Check 6
\draw[arrow, green!50!black] (gate.west) -- (sentinel.east)
    node[midway, above, font=\scriptsize, text=green!50!black] {Check 6};

% Bottom: Delivered to User
\node[box] (user) at (0, -8.8) {\textbf{Delivered to User}};
\draw[arrow] (verifier) -- (user);

% Sentinel JSON passed to Layer 2
\draw[arrow, blue!50!black, dashed] ([xshift=1.2cm]sentinel.south) -- ([xshift=1.2cm]verifier.north)
    node[midway, right, font=\scriptsize, text=blue!50!black] {sentinel JSON};

\end{tikzpicture}
\caption{Completion Integrity System architecture. Model output passes through three verification layers before delivery. Layer~1 (Completion Sentinel) fires as a stop hook, parsing the session transcript and running six mechanical checks. If a check fails, the message is blocked and feedback is returned to the model via stderr. Layer~2 (File Verifier) reads Layer~1 findings and rechecks against the filesystem. Layer~3 (VIBE\_GATE markers) operates during skill execution: values are produced by shell command substitution, not by model generation, and feed into Layer~1's Check~6.}
\label{fig:architecture}
\end{figure}
